\documentclass[UTF8]{ctexart}
\usepackage{amsmath,amssymb}
\usepackage[a4paper,margin=2.5cm]{geometry}
\title{几何与代数 1008}
\date{}

\begin{document}
\maketitle

\subsection*{二次曲面}

\begin{itemize}
\item \textbf{§1 球面}\quad 到定点的距离等于定长的点的集合。
\item \textbf{§2 柱面}\quad 准线 $+$ 母线方向。

\textbf{圆柱面}：准线是圆，母线方向和圆垂直。

\textbf{方程特点}：若母线方向是 $z$ 轴，则柱面方程是
$F(x,y)=0$。
\item \textbf{§3 锥面}\quad 顶点 $+$ 准线。

\textbf{圆锥面}：准线是圆，顶点在过圆心且垂直于准线的直线上。
\end{itemize}

\textbf{锥面方程的特点}：若顶点是原点，则锥面方程是齐次方程，即
\[
F(x,y,z)=0\ \Longrightarrow\ \forall t,\ F(tx,ty,tz)=0 .
\]
\textbf{例如}
\[
x^2+2pyz=0 .
\]

\subsection*{1.\ 在 $\mathbb R^2$ 上：$xy=1$ 是双曲线}

用坐标变换消去交叉项：
\[
\begin{cases}
\vec e_1'=\dfrac{\sqrt2}{2}\vec e_1+\dfrac{\sqrt2}{2}\vec e_2,\\[8pt]
\vec e_2'=-\dfrac{\sqrt2}{2}\vec e_1+\dfrac{\sqrt2}{2}\vec e_2,
\end{cases}
\]
这是把坐标轴绕原点逆时针旋转 $\dfrac\pi4$，对应
\[
x=\frac{x'+y'}{\sqrt2},\qquad
y=\frac{y'-x'}{\sqrt2}.
\]
代入 $xy=1$：
\[
\frac{x'+y'}{\sqrt2}\cdot\frac{y'-x'}{\sqrt2}=1
\ \Longrightarrow\
\frac{y'^2-x'^2}{2}=1
\ \Longrightarrow\
\frac{y'^2}{2}-\frac{x'^2}{2}=1 .
\]
这是双曲线的标准形式，两个半轴长都是 $\sqrt2$。

\subsection*{2.\ 在 $\mathbb R^3$ 上：$xy+yz+xz=0$ 消去交叉项}

取一组互相垂直的单位基向量（其中 $\vec e_2'$ 由 $\vec e_3'$ 与
$\vec e_1'$ 叉乘得到）：
\[
\begin{cases}
\vec e_1'=\dfrac{1}{\sqrt2}(1,-1,0),\\[8pt]
\vec e_2'=\vec e_3'\times\vec e_1'=\dfrac{1}{\sqrt6}(1,1,-2),\\[8pt]
\vec e_3'=\dfrac{1}{\sqrt3}(1,1,1).
\end{cases}
\]
由坐标变换公式，把旧坐标用新坐标 $(x',y',z')$ 表示：
\[
\begin{cases}
x=\dfrac{x'}{\sqrt2}+\dfrac{y'}{\sqrt6}+\dfrac{z'}{\sqrt3},\\[8pt]
y=-\dfrac{x'}{\sqrt2}+\dfrac{y'}{\sqrt6}+\dfrac{z'}{\sqrt3},\\[8pt]
z=-\dfrac{2y'}{\sqrt6}+\dfrac{z'}{\sqrt3}.
\end{cases}
\]

代入 $xy+yz+xz$，三项分别展开：
\[
\begin{aligned}
xy&=\left(\frac{x'}{\sqrt2}+\frac{y'}{\sqrt6}+\frac{z'}{\sqrt3}\right)
\left(-\frac{x'}{\sqrt2}+\frac{y'}{\sqrt6}+\frac{z'}{\sqrt3}\right)\\
&=-\frac{x'^2}{2}+\frac{y'^2}{6}+\frac{z'^2}{3}
-\frac{x'y'}{\sqrt3}+\frac{x'z'}{\sqrt6}-\frac{\sqrt2}{3}y'z',
\end{aligned}
\]
\[
\begin{aligned}
yz&=\left(-\frac{2y'}{\sqrt6}+\frac{z'}{\sqrt3}\right)
\left(-\frac{2y'}{\sqrt6}+\frac{z'}{\sqrt3}\right)
=\frac{2y'^2}{3}-\frac{\sqrt2}{3}y'z'+\frac{z'^2}{3},
\end{aligned}
\]
\[
\begin{aligned}
xz&=\left(\frac{x'}{\sqrt2}+\frac{y'}{\sqrt6}+\frac{z'}{\sqrt3}\right)
\left(-\frac{2y'}{\sqrt6}+\frac{z'}{\sqrt3}\right)\\
&=\frac{x'y'}{\sqrt3}+\frac{x'z'}{\sqrt6}-\frac{y'^2}{3}+\frac{\sqrt2}{3}y'z'+\frac{z'^2}{3}.
\end{aligned}
\]

带入可知
\[
xy+yz+xz=-\frac{x'^2}{2}-\frac{y'^2}{2}+z'^2 .
\]

于是 $xy+yz+xz=0$ 化为
\[
\frac{z'^2}{1}=\frac{x'^2}{2}+\frac{y'^2}{2},
\qquad\text{即}\qquad
2z'^2=x'^2+y'^2 .
\]
这是顶点在原点、轴沿 $z'$ 方向的\textbf{圆锥面}，且是齐次方程，
与前面“锥面方程的特点”一致。

\end{document}
